\documentclass[11pt,letterpaper]{article}

\usepackage[utf8]{inputenc}
\usepackage[T1]{fontenc}
\usepackage{lmodern}
\usepackage{microtype}
\usepackage{geometry}
\usepackage{amsmath,amssymb,amsthm,mathtools}
\usepackage{booktabs,tabularx,array}
\usepackage{float}
\usepackage{enumitem}
\usepackage{xcolor}
\usepackage{listings}
\usepackage{fancyhdr}
\usepackage{hyperref}
\usepackage{bookmark}

\geometry{margin=0.90in,headheight=14pt}
\allowdisplaybreaks[1]
\setlength{\parindent}{1.25em}
\setlength{\parskip}{0.25em}
\setlist[itemize]{leftmargin=1.6em,itemsep=0.2em,topsep=0.3em}
\emergencystretch=2em

\definecolor{codebg}{RGB}{246,247,249}
\definecolor{codeframe}{RGB}{210,214,220}
\lstdefinestyle{sage}{
  backgroundcolor=\color{codebg},
  frame=single,
  rulecolor=\color{codeframe},
  basicstyle=\ttfamily\scriptsize,
  columns=fullflexible,
  keepspaces=true,
  showstringspaces=false,
  breaklines=true,
  xleftmargin=0.35em,
  xrightmargin=0.35em,
  aboveskip=0.7em,
  belowskip=0.7em
}
\lstdefinestyle{output}{
  style=sage,
  language={}
}

\newtheorem{theorem}{Teorema}
\newtheorem{corollary}{Corolário}
\theoremstyle{definition}
\newtheorem{definition}{Definição}

\newcommand{\Dd}[2]{D_{#1}\!\left(#2\right)}
\newcommand{\Layer}[1]{\mathcal{L}_{#1}}
\newcommand{\PhiH}[3]{\Phi_{#1}\!\left(#2,#3\right)}
\renewcommand{\tablename}{Tabela}
\renewcommand{\refname}{Referências}

\pagestyle{fancy}
\fancyhf{}
\fancyhead[C]{\small Interseção entre diferenças progressivas cúbicas e quárticas}
\fancyfoot[C]{\thepage}
\renewcommand{\headrulewidth}{0pt}
\fancypagestyle{plain}{
  \fancyhf{}
  \fancyfoot[C]{\thepage}
  \renewcommand{\headrulewidth}{0pt}
}

\hypersetup{
  pdftitle={Interseção entre diferenças progressivas cúbicas e quárticas},
  pdfauthor={Zaqueu Ribeiro da Costa},
  pdfsubject={Classificação completa de D3(n)=D4(m) por meio de uma curva elíptica},
  pdfkeywords={diferenças progressivas, curvas elípticas, pontos inteiros, formas ciclotômicas},
  colorlinks=true,
  linkcolor=black,
  citecolor=black,
  urlcolor=blue
}

\title{\textbf{Interseção entre diferenças progressivas cúbicas e quárticas}\\[0.45em]
\large\itshape Classificação completa por meio de uma curva elíptica}
\author{\textbf{Zaqueu Ribeiro da Costa}\\
\small Pesquisador independente, Almenara, Minas Gerais, Brasil}
\date{\small Pré-publicação, versão 1, 17 de setembro de 2026}

\begin{document}

\maketitle

\renewcommand{\abstractname}{Resumo}
\begin{abstract}
Definimos
\[
  \Dd{d}{n}=(n+1)^d-n^d
\]
e determinamos todas as soluções inteiras não negativas da igualdade
\(\Dd{3}{n}=\Dd{4}{m}\).
Uma mudança de variáveis envia cada solução a um ponto inteiro da curva elíptica
\[
  E:\quad Y^2=X^3+36X-108
\]
sujeito a duas congruências. O cálculo certificado do grupo de Mordell-Weil fornece posto \(1\), torção trivial e gerador \((6,18)\). A enumeração completa dos pontos inteiros de \(E\) produz a classificação final. O resultado principal é que as únicas soluções não negativas são \((0,0)\) e \((70,15)\), e a única interseção positiva das camadas é \(14911\).

\medskip
\noindent\textbf{Palavras-chave:} diferenças progressivas; equações diofantinas; curvas elípticas; pontos inteiros; formas ciclotômicas.
\end{abstract}

\renewcommand{\abstractname}{Abstract}
\begin{abstract}
Let
\[
  \Dd{d}{n}=(n+1)^d-n^d.
\]
We determine all nonnegative integer solutions of
\(\Dd{3}{n}=\Dd{4}{m}\).
An exact change of variables maps every solution to an integral point on the elliptic curve
\[
  E:\quad Y^2=X^3+36X-108
\]
satisfying two congruence conditions. A proof-enabled computation of the Mordell-Weil group gives rank \(1\), trivial torsion, and generator \((6,18)\). The complete integral-point set of \(E\) yields the final classification. The main result is that the only nonnegative solutions are \((0,0)\) and \((70,15)\), and the only positive common value of the layers is \(14911\).

\medskip
\noindent\textbf{Keywords:} forward differences; Diophantine equations; elliptic curves; integral points; cyclotomic forms.
\end{abstract}

\section{Introdução}

A primeira diferença progressiva da função potência transforma uma sequência de potências em uma família de polinômios inteiros. Interseções entre os conjuntos de valores desses polinômios são o tema do presente trabalho.

\begin{equation}
  \Dd{d}{n}=(n+1)^d-n^d=\Delta(n^d).
  \label{eq:forward}
\end{equation}

O problema considerado é determinar os pares de inteiros positivos \((n,m)\) para os quais
\(\Dd{3}{n}=\Dd{4}{m}\).
Uma busca finita identifica a igualdade
\(\Dd{3}{70}=\Dd{4}{15}=14911\),
mas não pode excluir soluções posteriores. A completude exige substituir o limite de busca por uma classificação de pontos integrais sobre uma curva elíptica.

O resultado central mostra que \(14911\) é o único valor positivo compartilhado. A demonstração tem duas partes. Primeiro, uma transformação algébrica exata converte a equação original em uma equação de curva elíptica. Em seguida, a enumeração completa dos pontos inteiros da curva e o filtro congruencial eliminam todas as soluções restantes.

\section{Resultado principal}

\begin{theorem}
Para \(n,m\in\mathbb{Z}_{\geq0}\), a igualdade
\[
  (n+1)^3-n^3=(m+1)^4-m^4
\]
ocorre se, e somente se,
\((n,m)=(0,0)\) ou \((n,m)=(70,15)\).
\end{theorem}

\begin{corollary}
Para as camadas positivas definidas em \eqref{eq:layer}, a única interseção entre os expoentes \(3\) e \(4\) é o valor \(14911\).
\end{corollary}

\begin{equation}
  \Layer{3}\cap\Layer{4}=\{14911\}.
  \label{eq:intersection}
\end{equation}

A solução positiva verifica-se diretamente:
\[
  \Dd{3}{70}=71^3-70^3=14911,
  \qquad
  \Dd{4}{15}=16^4-15^4=14911.
\]

\section{Redução a uma curva elíptica}

A expansão das duas diferenças fornece
\begin{equation}
  \Dd{3}{n}=3n^2+3n+1,
  \qquad
  \Dd{4}{m}=4m^3+6m^2+4m+1.
  \label{eq:expansion}
\end{equation}

Introduza as variáveis ímpares
\(u=2m+1\) e \(v=2n+1\). Então
\begin{equation}
  \Dd{4}{m}=\frac{u^3+u}{2},
  \qquad
  \Dd{3}{n}=\frac{3v^2+1}{4}.
  \label{eq:uv}
\end{equation}

Logo, \(\Dd{3}{n}=\Dd{4}{m}\) equivale exatamente a
\begin{equation}
  3v^2=2u^3+2u-1.
  \label{eq:uvcurve}
\end{equation}

Escalando as coordenadas por
\begin{equation}
  X=6u=6(2m+1),
  \qquad
  Y=18v=18(2n+1),
  \label{eq:coordinates}
\end{equation}
obtém-se a equação de Weierstrass
\begin{equation}
  E:\qquad Y^2=X^3+36X-108.
  \label{eq:curve}
\end{equation}

A curva é não singular, pois seu discriminante é diferente de zero:
\begin{equation}
  \Delta_E
  =-16\left(4\cdot36^3+27\cdot(-108)^2\right)
  =-8024832\neq0.
  \label{eq:discriminant}
\end{equation}

\subsection{Filtro congruencial}

A transformação não envia as soluções para pontos inteiros arbitrários. Como \(u\) e \(v\) são ímpares positivos quando \(m,n\geq0\), suas imagens satisfazem
\begin{equation}
  X>0,\qquad Y>0,\qquad
  X\equiv6\pmod{12},\qquad
  Y\equiv18\pmod{36}.
  \label{eq:filter}
\end{equation}

Reciprocamente, qualquer ponto inteiro de \(E\) que satisfaça \eqref{eq:filter} determina
\begin{equation}
  m=\frac{X-6}{12},
  \qquad
  n=\frac{Y-18}{36}.
  \label{eq:inverse}
\end{equation}

\section{Certificado do grupo de Mordell-Weil e dos pontos inteiros}

O cálculo com SageMath, executado com \texttt{proof.all(True)}, produz a curva
\(Y^2=X^3+36X-108\), com posto \(1\), torção trivial e gerador \((6,18)\). A rotina de pontos inteiros fornece exatamente
\begin{equation}
  E(\mathbb{Z})=
  \{(4,\pm10),(6,\pm18),(186,\pm2538)\}.
\end{equation}

Aplicando o filtro congruencial \(X\equiv6\pmod{12}\) e \(Y\equiv18\pmod{36}\), só sobrevivem
\((6,18)\) e \((186,2538)\), equivalentes a \((n,m)=(0,0)\) e \((70,15)\).

Portanto, as únicas soluções não negativas são \
\((0,0)\) e \((70,15)\), e a única interseção positiva das camadas é \(14911\).

\section{Conclusão}

A igualdade entre a primeira diferença cúbica e a primeira diferença quártica reduz-se à curva elíptica
\(Y^2=X^3+36X-108\).
A determinação completa de \(E(\mathbb{Z})\), combinada com o filtro
\(X\equiv6\pmod{12}\) e \(Y\equiv18\pmod{36}\),
prova que a única colisão positiva é
\[
  \Dd{3}{70}=\Dd{4}{15}=14911.
\]

\textbf{Autor: Zaqueu Ribeiro da Costa} \
\textbf{Pesquisador independente}

\end{document}
